% Generated draft: requires mathematical and source review.
\documentclass[11pt]{article}
\usepackage[margin=25mm]{geometry}
\usepackage{fontspec}
\setmainfont{Latin Modern Roman}
\usepackage{amsmath,amssymb,mathtools,unicode-math}
\setmathfont{Latin Modern Math}
\usepackage{enumitem,needspace,xurl,hyperref}
\setlength{\parindent}{0pt}
\setlength{\parskip}{5pt}
\setlength{\emergencystretch}{3em}
\newcommand{\N}{\mathbb{N}}
\newcommand{\Z}{\mathbb{Z}}
\newcommand{\Q}{\mathbb{Q}}
\newcommand{\R}{\mathbb{R}}
\newcommand{\C}{\mathbb{C}}
\newcommand{\F}{\mathbb{F}}
\newcommand{\A}{\mathbb{A}}
\newcommand{\PP}{\mathbb{P}}
\newcommand{\E}{\mathbb{E}}
\newcommand{\eps}{\varepsilon}
\newcommand{\dd}{\,\mathrm{d}}
\begin{document}
\section*{UnsolvedMath ID: 6200027}
\textbf{Source problem code:} catalogue entry
\subsection{Definitions and mathematical statement}
A finitely generated group is two-ended if its Cayley graph has exactly two ends; equivalently it is virtually infinite cyclic. Splitting over a two-ended subgroup means a nontrivial amalgamated-product or HNN decomposition with that edge group, interpreted through Bass--Serre theory.

The source asks whether existence of such a splitting is an invariant of Bestvina boundaries. Under the topological-boundary interpretation, the proposed assertion is: if groups $G,H$ admit Bestvina boundaries $Z_G,Z_H$ and $Z_G\cong Z_H$, then
\[
G\text{ splits nontrivially over a two-ended subgroup}
\quad\Longleftrightarrow\quad
H\text{ splits nontrivially over a two-ended subgroup}.
\]
Bestvina boundaries ordinarily come from $Z$-compactifications associated with geometric group actions. The record gives neither its precise $Z$-structure definition nor the allowed group/action category. Homeomorphism is the displayed candidate equivalence, not an inferred assertion about shape equivalence or an equivariant boundary map.

The record appends three necessary conditions for a compact metrizable space $X$ to be a boundary of a proper cocompact CAT(0) space: it has one, two, or infinitely many components; it is finite dimensional (attributed to Swenson); and it has nontrivial top-degree \v{C}ech cohomology (attributed to Geoghegan--Ontaneda). These are supplied contextual conditions, not asserted sufficient conditions. Coefficients, reduced versus unreduced cohomology, and the exact space/action hypotheses are not specified.

The source then breaks off at ``In the case when $X$ admits a cocompact free(?''. No additional hypothesis or conclusion is reconstructed from that incomplete fragment.
\subsection{Short English statement}
Does the topology of a Bestvina boundary determine whether its group splits over a two-ended subgroup? The source also lists necessary CAT(0)-boundary conditions but ends in an incomplete sentence.
\subsection{Sources}
\begin{itemize}
\item Local UnsolvedMath registry, permanent ID 6200027, exact source code catalogue entry, ``Boundaries of Groups and Kleinian Groups---Problem 27''; proposed by Bruce Kleiner, 2005. Registry status open, classification OPEN-TRIAGE.
\item Kapovich, \emph{Problems on Boundaries of Groups and Kleinian Groups} (2005), Problem 27, PDF page 8, \url{https://www.math.ucdavis.edu/~kapovich/EPR/problems.pdf}. Supplied normalized pdf-text extraction accessed July 29, 2026. The retained statement is visibly truncated; no PDF was accessed here.
\item The source names Swenson for finite-dimensional CAT(0) boundaries and Geoghegan--Ontaneda for top \v{C}ech cohomology, without precise citations. The research summary names Papasoglu and Bowditch for related quasi-isometry invariance in hyperbolic groups, without theorem hypotheses or publication details, and says the Bestvina/CAT(0) boundary formulation remains open. No unrestricted transfer from hyperbolic to general CAT(0) groups is asserted.
\item Provenance flags current status and statement-reuse terms NEEDS\_REVIEW; the source is publicly author-hosted.
\end{itemize}
\end{document}
